\IfFileExists{stacks-project.cls}{%
\documentclass{stacks-project}
}{%
\documentclass{amsart}
}

% For dealing with references we use the comment environment
\usepackage{verbatim}
\newenvironment{reference}{\comment}{\endcomment}
%\newenvironment{reference}{}{}
\newenvironment{slogan}{\comment}{\endcomment}
\newenvironment{history}{\comment}{\endcomment}

% For commutative diagrams we use Xy-pic
\usepackage[all]{xy}

% We use 2cell for 2-commutative diagrams.
\xyoption{2cell}
\UseAllTwocells

% We use multicol for the list of chapters between chapters
\usepackage{multicol}

% This is generally recommended for better output
\usepackage{lmodern}
\usepackage[T1]{fontenc}

% For cross-file-references
\usepackage{xr-hyper}

% Package for hypertext links:
\usepackage{hyperref}

% For any local file, say "hello.tex" you want to link to please
% use \externaldocument[hello-]{hello}
\externaldocument[introduction-]{introduction}
\externaldocument[conventions-]{conventions}
\externaldocument[sets-]{sets}
\externaldocument[categories-]{categories}
\externaldocument[topology-]{topology}
\externaldocument[sheaves-]{sheaves}
\externaldocument[sites-]{sites}
\externaldocument[stacks-]{stacks}
\externaldocument[fields-]{fields}
\externaldocument[algebra-]{algebra}
\externaldocument[brauer-]{brauer}
\externaldocument[homology-]{homology}
\externaldocument[derived-]{derived}
\externaldocument[simplicial-]{simplicial}
\externaldocument[more-algebra-]{more-algebra}
\externaldocument[smoothing-]{smoothing}
\externaldocument[modules-]{modules}
\externaldocument[sites-modules-]{sites-modules}
\externaldocument[injectives-]{injectives}
\externaldocument[cohomology-]{cohomology}
\externaldocument[sites-cohomology-]{sites-cohomology}
\externaldocument[dga-]{dga}
\externaldocument[dpa-]{dpa}
\externaldocument[sdga-]{sdga}
\externaldocument[hypercovering-]{hypercovering}
\externaldocument[schemes-]{schemes}
\externaldocument[constructions-]{constructions}
\externaldocument[properties-]{properties}
\externaldocument[morphisms-]{morphisms}
\externaldocument[coherent-]{coherent}
\externaldocument[divisors-]{divisors}
\externaldocument[limits-]{limits}
\externaldocument[varieties-]{varieties}
\externaldocument[topologies-]{topologies}
\externaldocument[descent-]{descent}
\externaldocument[perfect-]{perfect}
\externaldocument[more-morphisms-]{more-morphisms}
\externaldocument[flat-]{flat}
\externaldocument[groupoids-]{groupoids}
\externaldocument[more-groupoids-]{more-groupoids}
\externaldocument[etale-]{etale}
\externaldocument[chow-]{chow}
\externaldocument[intersection-]{intersection}
\externaldocument[pic-]{pic}
\externaldocument[weil-]{weil}
\externaldocument[adequate-]{adequate}
\externaldocument[dualizing-]{dualizing}
\externaldocument[duality-]{duality}
\externaldocument[discriminant-]{discriminant}
\externaldocument[derham-]{derham}
\externaldocument[local-cohomology-]{local-cohomology}
\externaldocument[algebraization-]{algebraization}
\externaldocument[curves-]{curves}
\externaldocument[resolve-]{resolve}
\externaldocument[models-]{models}
\externaldocument[functors-]{functors}
\externaldocument[equiv-]{equiv}
\externaldocument[pione-]{pione}
\externaldocument[etale-cohomology-]{etale-cohomology}
\externaldocument[proetale-]{proetale}
\externaldocument[relative-cycles-]{relative-cycles}
\externaldocument[more-etale-]{more-etale}
\externaldocument[trace-]{trace}
\externaldocument[crystalline-]{crystalline}
\externaldocument[spaces-]{spaces}
\externaldocument[spaces-properties-]{spaces-properties}
\externaldocument[spaces-morphisms-]{spaces-morphisms}
\externaldocument[decent-spaces-]{decent-spaces}
\externaldocument[spaces-cohomology-]{spaces-cohomology}
\externaldocument[spaces-limits-]{spaces-limits}
\externaldocument[spaces-divisors-]{spaces-divisors}
\externaldocument[spaces-over-fields-]{spaces-over-fields}
\externaldocument[spaces-topologies-]{spaces-topologies}
\externaldocument[spaces-descent-]{spaces-descent}
\externaldocument[spaces-perfect-]{spaces-perfect}
\externaldocument[spaces-more-morphisms-]{spaces-more-morphisms}
\externaldocument[spaces-flat-]{spaces-flat}
\externaldocument[spaces-groupoids-]{spaces-groupoids}
\externaldocument[spaces-more-groupoids-]{spaces-more-groupoids}
\externaldocument[bootstrap-]{bootstrap}
\externaldocument[spaces-pushouts-]{spaces-pushouts}
\externaldocument[spaces-chow-]{spaces-chow}
\externaldocument[groupoids-quotients-]{groupoids-quotients}
\externaldocument[spaces-more-cohomology-]{spaces-more-cohomology}
\externaldocument[spaces-simplicial-]{spaces-simplicial}
\externaldocument[spaces-duality-]{spaces-duality}
\externaldocument[formal-spaces-]{formal-spaces}
\externaldocument[restricted-]{restricted}
\externaldocument[spaces-resolve-]{spaces-resolve}
\externaldocument[formal-defos-]{formal-defos}
\externaldocument[defos-]{defos}
\externaldocument[cotangent-]{cotangent}
\externaldocument[examples-defos-]{examples-defos}
\externaldocument[algebraic-]{algebraic}
\externaldocument[examples-stacks-]{examples-stacks}
\externaldocument[stacks-sheaves-]{stacks-sheaves}
\externaldocument[criteria-]{criteria}
\externaldocument[artin-]{artin}
\externaldocument[quot-]{quot}
\externaldocument[stacks-properties-]{stacks-properties}
\externaldocument[stacks-morphisms-]{stacks-morphisms}
\externaldocument[stacks-limits-]{stacks-limits}
\externaldocument[stacks-cohomology-]{stacks-cohomology}
\externaldocument[stacks-perfect-]{stacks-perfect}
\externaldocument[stacks-introduction-]{stacks-introduction}
\externaldocument[stacks-more-morphisms-]{stacks-more-morphisms}
\externaldocument[stacks-geometry-]{stacks-geometry}
\externaldocument[moduli-]{moduli}
\externaldocument[moduli-curves-]{moduli-curves}
\externaldocument[examples-]{examples}
\externaldocument[exercises-]{exercises}
\externaldocument[guide-]{guide}
\externaldocument[desirables-]{desirables}
\externaldocument[coding-]{coding}
\externaldocument[obsolete-]{obsolete}
\externaldocument[fdl-]{fdl}
\externaldocument[index-]{index}

% Theorem environments.
%
\theoremstyle{plain}
\newtheorem{theorem}[subsection]{Theorem}
\newtheorem{proposition}[subsection]{Proposition}
\newtheorem{lemma}[subsection]{Lemma}

\theoremstyle{definition}
\newtheorem{definition}[subsection]{Definition}
\newtheorem{example}[subsection]{Example}
\newtheorem{exercise}[subsection]{Exercise}
\newtheorem{situation}[subsection]{Situation}

\theoremstyle{remark}
\newtheorem{remark}[subsection]{Remark}
\newtheorem{remarks}[subsection]{Remarks}

\numberwithin{equation}{subsection}

% Macros
%
\def\lim{\mathop{\mathrm{lim}}\nolimits}
\def\colim{\mathop{\mathrm{colim}}\nolimits}
\def\Spec{\mathop{\mathrm{Spec}}}
\def\Hom{\mathop{\mathrm{Hom}}\nolimits}
\def\Ext{\mathop{\mathrm{Ext}}\nolimits}
\def\SheafHom{\mathop{\mathcal{H}\!\mathit{om}}\nolimits}
\def\SheafExt{\mathop{\mathcal{E}\!\mathit{xt}}\nolimits}
\def\Sch{\mathit{Sch}}
\def\Mor{\mathop{\mathrm{Mor}}\nolimits}
\def\Ob{\mathop{\mathrm{Ob}}\nolimits}
\def\Sh{\mathop{\mathit{Sh}}\nolimits}
\def\NL{\mathop{N\!L}\nolimits}
\def\CH{\mathop{\mathrm{CH}}\nolimits}
\def\proetale{{pro\text{-}\acute{e}tale}}
\def\etale{{\acute{e}tale}}
\def\QCoh{\mathit{QCoh}}
\def\Ker{\mathop{\mathrm{Ker}}}
\def\Im{\mathop{\mathrm{Im}}}
\def\Coker{\mathop{\mathrm{Coker}}}
\def\Coim{\mathop{\mathrm{Coim}}}

% Boxtimes
%
\DeclareMathSymbol{\boxtimes}{\mathbin}{AMSa}{"02}

%
% Macros for moduli stacks/spaces
%
\def\QCohstack{\mathcal{QC}\!\mathit{oh}}
\def\Cohstack{\mathcal{C}\!\mathit{oh}}
\def\Spacesstack{\mathcal{S}\!\mathit{paces}}
\def\Quotfunctor{\mathrm{Quot}}
\def\Hilbfunctor{\mathrm{Hilb}}
\def\Curvesstack{\mathcal{C}\!\mathit{urves}}
\def\Polarizedstack{\mathcal{P}\!\mathit{olarized}}
\def\Complexesstack{\mathcal{C}\!\mathit{omplexes}}
% \Pic is the operator that assigns to X its picard group, usage \Pic(X)
% \Picardstack_{X/B} denotes the Picard stack of X over B
% \Picardfunctor_{X/B} denotes the Picard functor of X over B
\def\Pic{\mathop{\mathrm{Pic}}\nolimits}
\def\Picardstack{\mathcal{P}\!\mathit{ic}}
\def\Picardfunctor{\mathrm{Pic}}
\def\Deformationcategory{\mathcal{D}\!\mathit{ef}}
\title{The cone of an admissible monomorphism is homotopy equivalent to the quotient, and the triangle is isomorphic}
\author{The Stacks Project Authors}
\date{stacks-project:dga.tex:local-snapshot-20261012; GNU FDL 1.2 or later}
\begin{document}
\maketitle
\noindent\textit{Editorial note.} The counted statement is the author's own lemma environment labelled lemma-cone-homotopy (source lines 4325-4379); the proof environment begins two lines below it and is its own proof, closing at source line 4489.\par\medskip
\noindent\textit{Editorial note.} The statement and the proof invoke several author results of the same chapter by label (the standing situation, the null-homotopy of the identity of a cone, and the factorisation of a morphism); they are retained as context spans, and the generator expands that context transitively, so every reference resolves inside the assembled package and no source pointer is declared.\par\medskip
\noindent\textit{Editorial note.} The statement contains two displayed Xy-pic diagrams of triangles and one labelled equation; all are part of the retained source range and are reproduced as drawn, so no figure file is needed.\par\medskip
\noindent\textit{Editorial note.} The retained preamble is the chapter's own minimal preamble (sources/stacks-dga/preamble.tex, taken with its real line count), a verbatim excerpt of the native Xy-pic setup, the native theorem declarations with the mathematical macros, and the project's own no-op reference environments; no macro was invented or redefined.\par\medskip
\noindent\textit{Editorial scope.} One original result of the Stacks Project's Differential Graded Algebra chapter, counted once: the cone of an admissible monomorphism is homotopy equivalent to the quotient, and the associated triangles are isomorphic.\par\medskip
\noindent\textit{Allowed standard prerequisites (editorial).}\begin{itemize}
\item The standing situation (axioms A, B, C). A dg category with the distinguished-triangle axioms, in particular the admissible short exact sequence attached to a morphism by axiom (C) and the associated triangle (the author's Situation situation-ABC, retained as context).. Axiom (C) supplies the sequence x -> C(alpha) -> x[1] and the structure maps b, p, s, pi, sigma used throughout (source lines 4327-4349, 4382-4410).
\item Cone of a morphism and its triangle. The cone C(f) of a morphism with its inclusion, projection and the triangle it forms, and the notion of a morphism of triangles.. The cone is the object being compared and the conclusion is an isomorphism of triangles (source lines 4331-4349, 4362-4377).
\item Null homotopy of the identity of a cone. The identity morphism on the cone of an identity is homotopic to zero (the author's lemma lemma-id-cone-null, retained as context).. It is invoked at the end of part (1) to conclude the homotopy equivalence (source lines 4463-4466).
\item Factorisation of a morphism through an admissible monomorphism. Every morphism factors through an admissible monomorphism which extends to an admissible short exact sequence (the author's lemma lemma-factor).. Part (2) applies it to alpha and transfers the isomorphism of triangles (source lines 4350-4377).
\end{itemize}\medskip
\begin{situation}
\label{situation-ABC}
Here $R$ is a ring and $\mathcal{A}$ is a differential graded category
over $R$ having axioms (A), (B), and
\begin{enumerate}
\item[(C)] given an arrow $f : x \to y$ of degree $0$ with
$\text{d}(f) = 0$ there exists an admissible short exact sequence
$y \to c(f) \to x[1]$ in $\text{Comp}(\mathcal{A})$ such that the map
$x[1] \to y[1]$ of Lemma \ref{lemma-get-triangle} is equal to $f[1]$.
\end{enumerate}
\end{situation}

\begin{lemma}
\label{lemma-id-cone-null}
In Situation \ref{situation-ABC}
given any object $x$ of $\mathcal{A}$, and the cone $C(1_x)$ of the
identity morphism $1_x : x \to x$, the identity morphism on
$C(1_x)$ is homotopic to zero.
\end{lemma}

\begin{proof}
Consider the admissible short exact sequence given by axiom (C).
$$
\xymatrix{
x \ar@<0.5ex>[r]^a  &
C(1_x) \ar@<0.5ex>[l]^{\pi} \ar@<0.5ex>[r]^b &
x[1]\ar@<0.5ex>[l]^s
}
$$
Then by Lemma \ref{lemma-get-triangle}, identifying hom-sets under
shifting, we have $1_x=\pi d(s)=-d(\pi)s$ where $s$ is regarded as
a morphism in $\Hom_{\mathcal{A}}^{-1}(x,C(1_x))$. Therefore
$a=a\pi d(s)=d(s)$ using formula (5) of Lemma \ref{lemma-get-triangle},
and $b=-d(\pi)sb=-d(\pi)$ by formula (6) of Lemma \ref{lemma-get-triangle}.
Hence
$$
1_{C(1_x)} = a\pi + sb = d(s)\pi - sd(\pi) = d(s\pi)
$$
since $s$ is of degree $-1$.
\end{proof}

\begin{lemma}
\label{lemma-factor}
In Situation \ref{situation-ABC} let $\alpha : x \to y$
be a morphism in $\text{Comp}(\mathcal{A})$. Then there exists
a factorization in $\text{Comp}(\mathcal{A})$:
$$
\xymatrix{
x \ar[r]^{\tilde{\alpha}}  &
\tilde{y} \ar@<0.5ex>[r]^{\pi} &
y\ar@<0.5ex>[l]^s
}
$$
such that
\begin{enumerate}
\item $\tilde{\alpha}$ is an admissible monomorphism, and
$\pi\tilde{\alpha}=\alpha$.
\item There exists a morphism
$s:y\to\tilde{y}$ in $\text{Comp}(\mathcal{A})$
such that $\pi s=1_y$ and $s\pi$ is homotopic to $1_{\tilde{y}}$. 
\end{enumerate}
\end{lemma}

\begin{proof}
By axiom (A), we may let $\tilde{y}$ be the differential graded direct
sum of $y$ and $C(1_x)$, i.e., there exists a diagram
$$
\xymatrix@C=3pc{
y \ar@<0.5ex>[r]^s  &
y\oplus C(1_x) \ar@<0.5ex>[l]^{\pi} \ar@<0.5ex>[r]^{p} &
C(1_x)\ar@<0.5ex>[l]^t
}
$$
where all morphisms are of degree zero, and in
$\text{Comp}(\mathcal{A})$. Let $\tilde{y} = y \oplus C(1_x)$.
Then $1_{\tilde{y}} = s\pi + tp$. Consider now the diagram
$$
\xymatrix{
x \ar[r]^{\tilde{\alpha}}  &
\tilde{y} \ar@<0.5ex>[r]^{\pi} &
y\ar@<0.5ex>[l]^s
}
$$
where $\tilde{\alpha}$ is induced by the morphism $x\xrightarrow{\alpha}y$
and the natural morphism $x\to C(1_x)$ fitting in the admissible
short exact sequence
$$
\xymatrix{
x \ar@<0.5ex>[r]  &
C(1_x) \ar@<0.5ex>[l] \ar@<0.5ex>[r] &
x[1]\ar@<0.5ex>[l]
}
$$
So the morphism $C(1_x)\to x$ of degree 0 in this diagram,
together with the zero morphism $y\to x$, induces a degree-0
morphism $\beta : \tilde{y} \to x$. Then $\tilde{\alpha}$ is an
admissible monomorphism since it fits into the admissible short
exact sequence
$$
\xymatrix{
x\ar[r]^{\tilde{\alpha}} &
\tilde{y} \ar[r] &
x[1]
}
$$
Furthermore, $\pi\tilde{\alpha} = \alpha$ by the construction of
$\tilde{\alpha}$, and $\pi s = 1_y$ by the first diagram. It
remains to show that $s\pi$ is homotopic to $1_{\tilde{y}}$.
Write $1_x$ as $d(h)$ for some degree $-1$ map. Then, our
last statement follows from
\begin{align*}
1_{\tilde{y}} - s\pi
= &
tp \\
= &
t(dh)p\quad\text{(by Lemma \ref{lemma-id-cone-null})} \\
= &
d(thp)
\end{align*}
since $dt = dp = 0$, and $t$ is of degree zero. 
\end{proof}

\begin{lemma}
\label{lemma-analogue-third-isomorphism}
In Situation \ref{situation-ABC} let $f_1 : x_1 \to y_1$ and
$f_2 : x_2 \to y_2$ be morphisms in $\text{Comp}(\mathcal{A})$. Let 
$$
(a,b,c): (x_1,y_1,c(f_1), f_1, i_1, p_1) \to (x_2,y_2, c(f_2), f_2, i_1, p_1)
$$
be any morphism of triangles in $K(\mathcal{A})$.
If $a$ and $b$ are homotopy equivalences, then so is $c$. 
\end{lemma}

\begin{proof}
Since $a$ and $b$ are homotopy equivalences, they are invertible in
$K(\mathcal{A})$ so let $a^{-1}$ and $b^{-1}$ denote their inverses
in $K(\mathcal{A})$, giving us a commutative diagram 
$$
\xymatrix{
x_2\ar[d]^{a^{-1}}\ar[r]^{f_2} &
y_2\ar[d]^{b^{-1}}\ar[r]^{i_2} &
c(f_2)\ar[d]^{c'} \\
x_1\ar[r]^{f_1} &
y_1 \ar[r]^{i_1} &
c(f_1)
}
$$
where the map $c'$ is defined via Lemma \ref{lemma-cone} applied to the left
commutative box of the above diagram. Since the diagram commutes
in $K(\mathcal{A})$, it suffices by Lemma \ref{lemma-triseq} to
prove the following: given a morphism of triangle
$(1,1,c): (x,y,c(f),f,i,p)\to (x,y,c(f),f,i,p)$
in $K(\mathcal{A})$, the map $c$ is an isomorphism in
$K(\mathcal{A})$. We have the commutative diagrams in $K(\mathcal{A})$:
$$
\vcenter{
\xymatrix{
y\ar[d]^{1}\ar[r] &
c(f)\ar[d]^{c}\ar[r] &
x[1]\ar[d]^{1} \\
y\ar[r] &
c(f) \ar[r] &
x[1]
}
}
\quad\Rightarrow\quad
\vcenter{
\xymatrix{
y\ar[d]^{0}\ar[r] &
c(f)\ar[d]^{c-1}\ar[r] &
x[1]\ar[d]^{0} \\
y\ar[r] &
c(f) \ar[r] &
x[1]
}
}
$$
Since the rows are admissible short exact sequences, we obtain
the identity $(c-1)^2 = 0$ by Lemma \ref{lemma-triseq}, from
which we conclude that $2-c$ is inverse to $c$ in $K(\mathcal{A})$
so that $c$ is an isomorphism. 
\end{proof}

\begin{lemma}
\label{lemma-cone}
\begin{slogan}
The homotopy category is a triangulated category.
This lemma proves a part of the axioms of a triangulated category.
\end{slogan}
In Situation \ref{situation-ABC} suppose that
$$
\xymatrix{
x_1 \ar[r]_{f_1} \ar[d]_a & y_1 \ar[d]^b \\
x_2 \ar[r]^{f_2} & y_2
}
$$
is a diagram of $\text{Comp}(\mathcal{A})$ commutative up to homotopy.
Then there exists a morphism $c : c(f_1) \to c(f_2)$ which gives rise to
a morphism of triangles
$$
(a, b, c) : (x_1, y_1, c(f_1)) \to (x_1, y_1, c(f_1))
$$
in $K(\mathcal{A})$.
\end{lemma}

\begin{proof}
The assumption means there exists a morphism $h : x_1 \to y_2$ of degree
$-1$ such that $\text{d}(h) = b f_1 - f_2 a$. Choose isomorphisms
$c(f_i) = y_i \oplus x_i[1]$ of graded objects compatible with the
morphisms $y_i \to c(f_i) \to x_i[1]$. Let's denote
$a_i : y_i \to c(f_i)$, $b_i : c(f_i) \to x_i[1]$, $s_i : x_i[1] \to c(f_i)$,
and $\pi_i : c(f_i) \to y_i$ the given morphisms. Recall that
$x_i[1] \to y_i[1]$ is given by $\pi_i \text{d}(s_i)$. By axiom (C)
this means that
$$
f_i = \pi_i \text{d}(s_i) = - \text{d}(\pi_i) s_i
$$
(we identify $\Hom(x_i, y_i)$ with $\Hom(x_i[1], y_i[1])$
using the shift functor $[1]$).
Set $c = a_2 b \pi_1 + s_2 a b_1 + a_2hb$. Then, using the
equalities found in the proof of Lemma \ref{lemma-get-triangle}
we obtain
\begin{align*}
\text{d}(c)
& =
a_2 b \text{d}(\pi_1) + \text{d}(s_2) a b_1 + a_2 \text{d}(h) b_1 \\
& =
- a_2 b f_1 b_1 + a_2 f_2 a b_1 + a_2 (b f_1 - f_2 a) b_1 \\
& = 0
\end{align*}
(where we have used in particular that
$\text{d}(\pi_1) = \text{d}(\pi_1) s_1 b_1 = f_1 b_1$ and
$\text{d}(s_2) = a_2 \pi_2 \text{d}(s_2) = a_2 f_2$).
Thus $c$ is a degree $0$ morphism $c : c(f_1) \to c(f_2)$ of $\mathcal{A}$
compatible with the given morphisms $y_i \to c(f_i) \to x_i[1]$.
\end{proof}

\begin{lemma}
\label{lemma-get-triangle}
Let $\mathcal{A}$ be a differential graded category satisfying
axioms (A) and (B). Given an admissible short exact sequence
$x \to y \to z$ we obtain (see proof) a triangle
$$
x \to y \to z \to x[1]
$$
in $\text{Comp}(\mathcal{A})$ with the property that any two compositions
in $z[-1] \to x \to y \to z \to x[1]$ are zero in $K(\mathcal{A})$.
\end{lemma}

\begin{proof}
Choose a diagram
$$
\xymatrix{
x \ar[rr]_1 \ar[rd]_a & & x \\
& y \ar[ru]_\pi \ar[rd]^b & \\
z \ar[rr]^1 \ar[ru]^s & & z
}
$$
giving the isomorphism of graded objects $y \cong x \oplus z$ as in the
definition of an admissible short exact sequence. Here are some equations
that hold in this situation
\begin{enumerate}
\item $1 = \pi a$ and hence $\text{d}(\pi) a = 0$,
\item $1 = b s$ and hence $b \text{d}(s) = 0$,
\item $1 = a \pi + s b$ and hence $a \text{d}(\pi) + \text{d}(s) b = 0$,
\item $\pi s = 0$ and hence $\text{d}(\pi)s + \pi \text{d}(s) = 0$,
\item $\text{d}(s) = a \pi \text{d}(s)$ because
$\text{d}(s) = (a \pi + s b)\text{d}(s)$ and $b\text{d}(s) = 0$,
\item $\text{d}(\pi) = \text{d}(\pi) s b$ because
$\text{d}(\pi) = \text{d}(\pi)(a \pi + s b)$ and $\text{d}(\pi)a = 0$,
\item $\text{d}(\pi \text{d}(s)) = 0$ because if we postcompose it
with the monomorphism $a$ we get
$\text{d}(a\pi \text{d}(s)) = \text{d}(\text{d}(s)) = 0$, and
\item $\text{d}(\text{d}(\pi)s) = 0$ as by (4) it is the negative
of $\text{d}(\pi\text{d}(s))$ which is $0$ by (7).
\end{enumerate}
We've used repeatedly that $\text{d}(a) = 0$, $\text{d}(b) = 0$,
and that $\text{d}(1) = 0$. By (7) we see that
$$
\delta = \pi \text{d}(s) = - \text{d}(\pi) s : z \to x[1]
$$
is a morphism in $\text{Comp}(\mathcal{A})$. By (5) we see that
the composition $a \delta = a \pi \text{d}(s) = \text{d}(s)$
is homotopic to zero. By (6) we see that the composition
$\delta b = - \text{d}(\pi)sb = \text{d}(-\pi)$ is homotopic to zero.
\end{proof}

\begin{lemma}
\label{lemma-triseq}
In Situation \ref{situation-ABC} let $x_i \to y_i \to z_i$
be morphisms in $\mathcal{A}$ ($i=1,2,3$) such that
$x_2 \to y_2\to z_2$ is an admissible short exact sequence.
Let $b : y_1 \to y_2$ and $b' : y_2\to y_3$ be morphisms
in $\text{Comp}(\mathcal{A})$ such that
$$
\vcenter{
\xymatrix{
x_1 \ar[d]_0 \ar[r] &
y_1 \ar[r] \ar[d]_b &
z_1 \ar[d]_0 \\
x_2 \ar[r] & y_2 \ar[r] & z_2
}
}
\quad\text{and}\quad
\vcenter{
\xymatrix{
x_2 \ar[d]^0 \ar[r] &
y_2 \ar[r] \ar[d]^{b'} &
z_2 \ar[d]^0 \\
x_3 \ar[r] & y_3 \ar[r] & z_3
}
}
$$
commute up to homotopy. Then $b'\circ b$ is homotopic to $0$.
\end{lemma}

\begin{proof}
By Lemma \ref{lemma-homo-change}, we can replace $b$ and $b'$
by homotopic maps $\tilde{b}$ and $\tilde{b}'$, such that the right
square of the left diagram commutes and the left square of the right
diagram commutes. Say $b = \tilde{b} + d(h)$ and $b'=\tilde{b}'+d(h')$
for degree $-1$ morphisms $h$ and $h'$ in $\mathcal{A}$. Hence
$$
b'b = \tilde{b}'\tilde{b} + d(\tilde{b}'h + h'\tilde{b} + h'd(h))
$$
since $d(\tilde{b})=d(\tilde{b}')=0$, i.e. $b'b$ is homotopic to
$\tilde{b}'\tilde{b}$. We now want to show that $\tilde{b}'\tilde{b}=0$.
Because $x_2\xrightarrow{f} y_2\xrightarrow{g} z_2$ is an admissible
short exact sequence, there exist degree $0$ morphisms
$\pi : y_2 \to x_2$ and $s : z_2 \to y_2$ such that
$\text{id}_{y_2} = f\pi + sg$. Therefore
$$
\tilde{b}'\tilde{b} = \tilde{b}'(f\pi + sg)\tilde{b} = 0
$$
since $g\tilde{b} = 0$ and $\tilde{b}'f = 0$ as consequences
of the two commuting squares.
\end{proof}

\begin{lemma}
\label{lemma-homo-change}
In Situation \ref{situation-ABC} given a diagram
$$
\xymatrix{x\ar[r]^f\ar[d]_a & y\ar[d]^b\\
z\ar[r]^g & w}
$$
in $\text{Comp}(\mathcal{A})$ commuting up to homotopy. Then
\begin{enumerate}
\item If $f$ is an admissible monomorphism, then $b$ is homotopic
to a morphism $b'$ which makes the diagram commute.
\item If $g$ is an admissible epimorphism, then $a$ is homotopic
to a morphism $a'$ which makes the diagram commute.
\end{enumerate}
\end{lemma}

\begin{proof}
To prove (1), observe that the hypothesis implies that there is some
$h\in\Hom_{\mathcal{A}}(x,w)$ of degree $-1$ such that $bf-ga=d(h)$.
Since $f$ is an admissible monomorphism, there is a morphism
$\pi : y \to x$ in the category $\mathcal{A}$ of degree $0$.
Let $b' = b - d(h\pi)$. Then
\begin{align*}
b'f = bf - d(h\pi)f
= &
bf - d(h\pi f) \quad (\text{since }d(f) = 0) \\
= &
bf-d(h) \\
= &
ga
\end{align*}
as desired. The proof for (2) is omitted.
\end{proof}

\begin{lemma}
\label{lemma-cone-homotopy}
In Situation \ref{situation-ABC}.
\begin{enumerate}
\item Given an admissible short exact sequence
$x\xrightarrow{\alpha} y\xrightarrow{\beta} z$.
Then there exists a homotopy equivalence
$e:C(\alpha)\to z$ such that the diagram
\begin{equation}
\label{equation-cone-isom-triangle}
\vcenter{
\xymatrix{
x\ar[r]^{\alpha}\ar[d] &
y\ar[r]^{b}\ar[d] &
C(\alpha)\ar[r]^{-c}\ar@{.>}[d]^{e} &
x[1]\ar[d] \\
x\ar[r]^{\alpha} &
y\ar[r]^{\beta} &
z\ar[r]^{\delta} & x[1]
}
}
\end{equation}
defines an isomorphism of triangles in $K(\mathcal{A})$. Here
$y\xrightarrow{b}C(\alpha)\xrightarrow{c}x[1]$
is the admissible short exact sequence given as in axiom (C).
\item Given a morphism
$\alpha : x \to y$ in $\text{Comp}(\mathcal{A})$, let
$x \xrightarrow{\tilde{\alpha}} \tilde{y} \to y$ be the
factorization given as in Lemma \ref{lemma-factor}, where the admissible
monomorphism $x \xrightarrow{\tilde{\alpha}} y$ extends to the
admissible short exact sequence
$$
\xymatrix{
x \ar[r]^{\tilde{\alpha}} &
\tilde{y} \ar[r] & z
}
$$
Then there exists an isomorphism of triangles
$$
\xymatrix{
x \ar[r]^{\tilde{\alpha}} \ar[d] &
\tilde{y} \ar[r] \ar[d] &
z \ar[r]^{\delta} \ar@{.>}[d]^{e} &
x[1] \ar[d] \\
x \ar[r]^{\alpha} &
y \ar[r] &
C(\alpha) \ar[r]^{-c} &
x[1]
}
$$
where the upper triangle is the triangle
associated to the sequence
$x \xrightarrow{\tilde{\alpha}} \tilde{y} \to z$.
\end{enumerate}
\end{lemma}

\begin{proof}
For (1), we consider the more complete diagram, \emph{without} the
sign change on $c$:
$$
\xymatrix{
x\ar@<0.5ex>[r]^{\alpha} \ar[d] &
y\ar@<0.5ex>[l]^{\pi} \ar@<0.5ex>[r]^{b}\ar[d] &
C(\alpha)\ar@<0.5ex>[l]^{p} \ar@<0.5ex>[r]^{c}\ar@{.>}@<0.5ex>[d]^{e} &
x[1]\ar@<0.5ex>[l]^{\sigma} \ar[d]\ar@<0.5ex>[r]^{\alpha} &
y[1]\ar@<0.5ex>[l]^{\pi} \\
x\ar@<0.5ex>[r]^{\alpha} &
y\ar@<0.5ex>[r]^{\beta} \ar@<0.5ex>[l]^{\pi} &
z\ar[r]^{\delta}\ar@<0.5ex>[l]^{s} \ar@{.>}@<0.5ex>[u]^{f} &
x[1]
}
$$
where the admissible short exact sequence
$x\xrightarrow{\alpha} y\xrightarrow{\beta} z$
is given the splitting $\pi$, $s$, and the admissible short exact sequence
$y\xrightarrow{b}C(\alpha)\xrightarrow{c}x[1]$ is given the splitting
$p$, $\sigma$. Note that (identifying hom-sets under shifting)
$$
\alpha = pd(\sigma) = -d(p)\sigma,\quad
\delta = \pi d(s) = -d(\pi)s
$$
by the construction in Lemma \ref{lemma-get-triangle}.

\medskip\noindent
We define $e=\beta p$ and $f=bs-\sigma\delta$. We first check that they are
morphisms in $\text{Comp}(\mathcal{A})$. To show that $d(e)=\beta d(p)$
vanishes, it suffices to show that $\beta d(p)b$ and $\beta d(p)\sigma$
both vanish, whereas
$$
\beta d(p)b = \beta d(pb) = \beta d(1_y) = 0,\quad
\beta d(p)\sigma = -\beta\alpha = 0
$$
Similarly, to check that $d(f)=bd(s)-d(\sigma)\delta$ vanishes,
it suffices to check the post-compositions by $p$ and $c$ both vanish,
whereas
\begin{align*}
pbd(s) - pd(\sigma)\delta
= &
d(s)-\alpha\delta = d(s)-\alpha\pi d(s) = 0 \\
cbd(s)-cd(\sigma)\delta
= &
-cd(\sigma)\delta = -d(c\sigma)\delta = 0
\end{align*}
The commutativity of left two squares of the
diagram \ref{equation-cone-isom-triangle} follows directly from definition.
Before we prove the commutativity of the right square (up to homotopy),
we first check that $e$ is a homotopy equivalence. Clearly,
$$
ef=\beta p (bs-\sigma\delta)=\beta s=1_z
$$
To check that $fe$ is homotopic to $1_{C(\alpha)}$, we first observe
$$
b\alpha = bpd(\alpha) = d(\sigma),\quad
\alpha c = -d(p)\sigma c = -d(p),\quad
d(\pi)p = d(\pi)s\beta p = -\delta\beta p
$$
Using these identities, we compute
\begin{align*}
1_{C(\alpha)} = &
bp + \sigma c
\quad (\text{from }y \xrightarrow{b} C(\alpha) \xrightarrow{c} x[1]) \\
= &
b(\alpha\pi + s\beta)p + \sigma(\pi\alpha)c
\quad (\text{from }x \xrightarrow{\alpha} y \xrightarrow{\beta} z) \\
= &
d(\sigma)\pi p + bs\beta p - \sigma\pi d(p)
\quad (\text{by the first two identities above}) \\
= &
d(\sigma)\pi p + bs\beta p - \sigma\delta\beta p
+ \sigma\delta\beta p - \sigma\pi d(p) \\
= &
(bs - \sigma\delta)\beta p + d(\sigma)\pi p
- \sigma d(\pi)p - \sigma\pi d(p)\quad
(\text{by the third identity above}) \\
= &
fe + d(\sigma \pi p)
\end{align*}
since $\sigma \in \Hom^{-1}(x, C(\alpha))$
(cf. proof of Lemma \ref{lemma-id-cone-null}).
Hence $e$ and $f$ are homotopy inverses.
Finally, to check that the right square of
diagram \ref{equation-cone-isom-triangle} commutes up to homotopy,
it suffices to check that $-cf=\delta$. This follows from
$$
-cf = -c(bs-\sigma\delta) = c\sigma\delta = \delta
$$
since $cb=0$.

\medskip\noindent
For (2), consider the factorization
$x\xrightarrow{\tilde{\alpha}}\tilde{y}\to y$
given as in Lemma \ref{lemma-factor}, so the second morphism
is a homotopy equivalence. By Lemmas \ref{lemma-cone} and
\ref{lemma-analogue-third-isomorphism}, there
exists an isomorphism of triangles between
$$
x \xrightarrow{\alpha} y \to C(\alpha) \to x[1]
\quad\text{and}\quad
x \xrightarrow{\tilde{\alpha}} \tilde{y} \to C(\tilde{\alpha}) \to x[1]
$$
Since we can compose isomorphisms of triangles, by replacing
$\alpha$ by $\tilde{\alpha}$, $y$ by $\tilde{y}$, and $C(\alpha)$ by
$C(\tilde{\alpha})$, we may assume $\alpha$ is an admissible monomorphism.
In this case, the result follows from (1).
\end{proof}

\end{document}
